\chapter{Introduction}

\projectname{} is a Unity-based tactical card battle prototype designed for teenagers and casual strategy players. The delivered software provides a complete single-battle session: the player enters through a start screen, controls allied units on a grid, plays tactical cards, observes enemy response, and reaches a clear result state.

\section{Motivation}

The project began from one simple design goal: each turn should ask the player to make a small number of visible and meaningful tactical decisions without requiring a large content set or a complicated ruleset. Instead of aiming for a campaign-heavy tactics game, the design focuses on one encounter in which positioning, range, timing, and action order matter.

The tactical card format was chosen because it combines two readable systems. First, a grid-based battle makes spatial decisions visible on screen. Second, a small hand of action cards introduces variety without turning the project into a deckbuilding or collection-heavy game. This makes the game strategically meaningful while still remaining approachable for the intended audience.

\section{Report Scope}

This report is written as a self-contained final software-engineering document rather than as a change log. The main chapters focus on the delivered system itself: background, requirements, development, validation, maintenance, and reflection. Supporting technical artifacts are preserved in the appendices, including the user manual, meeting minutes, representative test evidence, development planning notes, version-control summary, and document register.

That structure is deliberate. It allows the main discussion to stay focused on the product and the engineering decisions behind it, while still showing that the submission satisfies the wider expectations of a software-engineering project report: a concrete software artifact, updated system and UI design, traceable requirement evolution, implementation strategy, validation evidence, and reflective discussion of technical and team-level work \cite{courseSpecification}.

\section{Aims and Objectives}

The project has four main objectives.

\begin{enumerate}
    \item Deliver a complete playable battle loop in Unity, including player turn, enemy turn, and battle-end handling.
    \item Keep the gameplay narrow enough that the core interactions remain readable, stable, and testable.
    \item Build the system so that turn flow, card resolution, grid interaction, and runtime UI are understandable at code level.
    \item Improve interface clarity and feedback so that a first-time player can understand what to do without outside explanation.
\end{enumerate}

Success is therefore defined less by feature quantity than by the completeness of one polished vertical slice. A successful outcome means that a player can launch the game, start a battle, understand the interaction order, observe enemy response, finish the encounter, and then replay or return cleanly to the menu.

\section{Delivered Contribution}

The delivered build is intentionally modest in content, but complete in structure. It contains:

\begin{itemize}
    \item a full start-to-result session flow;
    \item a 6x6 tactical battle with card-driven actions and deterministic enemy response;
    \item a panel-based HUD with direct action feedback, prompts, and world-space health bars;
    \item automated validation for deterministic runtime behaviour, plus manual smoke for presentation-heavy UI;
    \item a code structure that separates flow, battle rules, UI, and presentation cleanly.
\end{itemize}

The project did not remain static from concept to completion. Early design assumptions were revised as playable builds exposed onboarding, readability, and session-flow problems. The report therefore does not present the final system as if it had emerged fully formed. Chapter 3 explains how the requirement set evolved from the earliest design baseline into the delivered implementation, and Chapter 4 explains how that evolution was handled through iterative development rather than a once-through sequence.
